\documentclass[10pt,a4paper]{amsart}
\usepackage{amsmath,amssymb,amsthm,graphicx}
\usepackage{verbatim}
\newenvironment{reference}{\comment}{\endcomment}
%\newenvironment{reference}{}{}
\newenvironment{slogan}{\comment}{\endcomment}
\newenvironment{history}{\comment}{\endcomment}

% For commutative diagrams we use Xy-pic
\usepackage[all]{xy}

% We use 2cell for 2-commutative diagrams.
\xyoption{2cell}
% Theorem environments.
%
\theoremstyle{plain}
\newtheorem{theorem}[subsection]{Theorem}
\newtheorem{proposition}[subsection]{Proposition}
\newtheorem{lemma}[subsection]{Lemma}

\theoremstyle{definition}
\newtheorem{definition}[subsection]{Definition}
\newtheorem{example}[subsection]{Example}
\newtheorem{exercise}[subsection]{Exercise}
\newtheorem{situation}[subsection]{Situation}

\theoremstyle{remark}
\newtheorem{remark}[subsection]{Remark}
\newtheorem{remarks}[subsection]{Remarks}

\numberwithin{equation}{subsection}

% Macros
%
\def\lim{\mathop{\mathrm{lim}}\nolimits}
\def\colim{\mathop{\mathrm{colim}}\nolimits}
\def\Spec{\mathop{\mathrm{Spec}}}
\def\Hom{\mathop{\mathrm{Hom}}\nolimits}
\def\Ext{\mathop{\mathrm{Ext}}\nolimits}
\def\SheafHom{\mathop{\mathcal{H}\!\mathit{om}}\nolimits}
\def\SheafExt{\mathop{\mathcal{E}\!\mathit{xt}}\nolimits}
\def\Sch{\mathit{Sch}}
\def\Mor{\mathop{\mathrm{Mor}}\nolimits}
\def\Ob{\mathop{\mathrm{Ob}}\nolimits}
\def\Sh{\mathop{\mathit{Sh}}\nolimits}
\def\NL{\mathop{N\!L}\nolimits}
\def\CH{\mathop{\mathrm{CH}}\nolimits}
\def\proetale{{pro\text{-}\acute{e}tale}}
\def\etale{{\acute{e}tale}}
\def\QCoh{\mathit{QCoh}}
\def\Ker{\mathop{\mathrm{Ker}}}
\def\Im{\mathop{\mathrm{Im}}}
\def\Coker{\mathop{\mathrm{Coker}}}
\def\Coim{\mathop{\mathrm{Coim}}}

% Boxtimes
%
\DeclareMathSymbol{\boxtimes}{\mathbin}{AMSa}{"02}

%
% Macros for moduli stacks/spaces
%
\def\QCohstack{\mathcal{QC}\!\mathit{oh}}
\def\Cohstack{\mathcal{C}\!\mathit{oh}}
\def\Spacesstack{\mathcal{S}\!\mathit{paces}}
\def\Quotfunctor{\mathrm{Quot}}
\def\Hilbfunctor{\mathrm{Hilb}}
\def\Curvesstack{\mathcal{C}\!\mathit{urves}}
\def\Polarizedstack{\mathcal{P}\!\mathit{olarized}}
\def\Complexesstack{\mathcal{C}\!\mathit{omplexes}}
% \Pic is the operator that assigns to X its picard group, usage \Pic(X)
% \Picardstack_{X/B} denotes the Picard stack of X over B
% \Picardfunctor_{X/B} denotes the Picard functor of X over B
\def\Pic{\mathop{\mathrm{Pic}}\nolimits}
\def\Picardstack{\mathcal{P}\!\mathit{ic}}
\def\Picardfunctor{\mathrm{Pic}}
\def\Deformationcategory{\mathcal{D}\!\mathit{ef}}
\usepackage{cleveref}
\crefname{originallemma}{original lemma}{original lemmas}
\Crefname{originallemma}{Original lemma}{Original lemmas}
\makeatletter
\newlabel{cohomology-lemma-tor-amplitude-stalk}{{09U9}{source}{}{sourcepointer-1}{}}
\newlabel{cohomology-lemma-tor-amplitude-stalk@cref}{{[originallemma][0][]09U9}{[1][0][]source}{}{}{}}
\makeatother
\crefname{originalequation}{original equation}{original equations}
\Crefname{originalequation}{Original equation}{Original equations}
\makeatletter
\newlabel{equation-epsilon}{{08H7}{source}{}{sourcepointer-2}{}}
\newlabel{equation-epsilon@cref}{{[originalequation][0][]08H7}{[1][0][]source}{}{}{}}
\makeatother
\crefname{originallemma}{original lemma}{original lemmas}
\Crefname{originallemma}{Original lemma}{Original lemmas}
\makeatletter
\newlabel{etale-cohomology-lemma-describe-etale-local-ring}{{04HX}{source}{}{sourcepointer-3}{}}
\newlabel{etale-cohomology-lemma-describe-etale-local-ring@cref}{{[originallemma][0][]04HX}{[1][0][]source}{}{}{}}
\makeatother
\crefname{originallemma}{original lemma}{original lemmas}
\Crefname{originallemma}{Original lemma}{Original lemmas}
\makeatletter
\newlabel{etale-cohomology-lemma-stalk-pullback}{{03Q1}{source}{}{sourcepointer-4}{}}
\newlabel{etale-cohomology-lemma-stalk-pullback@cref}{{[originallemma][0][]03Q1}{[1][0][]source}{}{}{}}
\makeatother
\crefname{originallemma}{original lemma}{original lemmas}
\Crefname{originallemma}{Original lemma}{Original lemmas}
\makeatletter
\newlabel{flat-lemma-tor-amplitude-up-down-henselization}{{0DK0}{source}{}{sourcepointer-5}{}}
\newlabel{flat-lemma-tor-amplitude-up-down-henselization@cref}{{[originallemma][0][]0DK0}{[1][0][]source}{}{}{}}
\makeatother
\crefname{originallemma}{original lemma}{original lemmas}
\Crefname{originallemma}{Original lemma}{Original lemmas}
\makeatletter
\newlabel{more-algebra-lemma-dumb-properties-henselization}{{07QM}{source}{}{sourcepointer-6}{}}
\newlabel{more-algebra-lemma-dumb-properties-henselization@cref}{{[originallemma][0][]07QM}{[1][0][]source}{}{}{}}
\makeatother
\crefname{originallemma}{original lemma}{original lemmas}
\Crefname{originallemma}{Original lemma}{Original lemmas}
\makeatletter
\newlabel{more-algebra-lemma-no-change-tor-amplitude}{{0DJF}{source}{}{sourcepointer-7}{}}
\newlabel{more-algebra-lemma-no-change-tor-amplitude@cref}{{[originallemma][0][]0DJF}{[1][0][]source}{}{}{}}
\makeatother
\crefname{originallemma}{original lemma}{original lemmas}
\Crefname{originallemma}{Original lemma}{Original lemmas}
\makeatletter
\newlabel{sheaves-lemma-stalk-pullback}{{008H}{source}{}{sourcepointer-8}{}}
\newlabel{sheaves-lemma-stalk-pullback@cref}{{[originallemma][0][]008H}{[1][0][]source}{}{}{}}
\makeatother
\crefname{originallemma}{original lemma}{original lemmas}
\Crefname{originallemma}{Original lemma}{Original lemmas}
\makeatletter
\newlabel{sites-cohomology-lemma-tor-amplitude-pullback}{{08H5}{source}{}{sourcepointer-9}{}}
\newlabel{sites-cohomology-lemma-tor-amplitude-pullback@cref}{{[originallemma][0][]08H5}{[1][0][]source}{}{}{}}
\makeatother
\title{Tor amplitude is unchanged by passing to the small etale site}
\author{The Stacks Project Authors}
\date{stacks-project:spaces-perfect.tex:local-snapshot-20261012; GNU FDL 1.2 or later}
\begin{document}
\maketitle
\noindent\textit{Editorial note.} All delivered mathematical text is the author's own, copied from the pinned Stacks Project chapter file spaces-perfect.tex; no statement or proof step was rewritten, reordered or removed.\par\medskip
\noindent\textit{Editorial note.} The counted proof is the author's complete argument: the proof is complete, covering both directions of (1) and deriving (2) formally\par\medskip
\noindent\textit{Editorial note.} The label lemma-tor-dimension-rel names the counted statement, and the proof environment immediately following it in the source is delivered verbatim as its proof.\par\medskip
\noindent\textit{Editorial note.} The counted statement is carried by the author's own lemma environment in the chapter source.\par\medskip
\noindent\textit{Allowed standard prerequisites (editorial).}\begin{itemize}
\item the localization functor epsilon\textasciicircum{}* from D(O\_X) to the small etale derived category over f\textasciicircum{}\{-1\}O\_Y
\item tor amplitude of complexes over a morphism of ringed spaces, checked on stalks
\item stalks of pullback sheaves at geometric points and strict henselization
\item faithful flatness of henselization and its invariance for tor amplitude
\end{itemize}\medskip
\ifdefined\hypertarget\hypertarget{sourcepointer-1}{}\fi
\noindent\textit{Source reference.} Original lemma 09U9: sources/stacks-chow/tags.txt, source lines 12417--12417. Prerequisite pointer invoked at proof line(s) 3281; it names a standard fact of a sibling chapter that the counted argument uses, and every step of the counted proof built on it is retained verbatim in the delivered body.\par\medskip
\ifdefined\hypertarget\hypertarget{sourcepointer-2}{}\fi
\noindent\textit{Source reference.} Original equation 08H7: sources/stacks-spaces-perfect/pinned-original-source.tex, source lines 320--320. Prerequisite pointer invoked at proof line(s) 3237; it names a standard fact of the same chapter that the counted argument uses, and every step of the counted proof built on it is retained verbatim in the delivered body.\par\medskip
\ifdefined\hypertarget\hypertarget{sourcepointer-3}{}\fi
\noindent\textit{Source reference.} Original lemma 04HX: sources/stacks-chow/tags.txt, source lines 5883--5883. Prerequisite pointer invoked at proof line(s) 3256; it names a standard fact of a sibling chapter that the counted argument uses, and every step of the counted proof built on it is retained verbatim in the delivered body.\par\medskip
\ifdefined\hypertarget\hypertarget{sourcepointer-4}{}\fi
\noindent\textit{Source reference.} Original lemma 03Q1: sources/stacks-chow/tags.txt, source lines 4878--4878. Prerequisite pointer invoked at proof line(s) 3255; it names a standard fact of a sibling chapter that the counted argument uses, and every step of the counted proof built on it is retained verbatim in the delivered body.\par\medskip
\ifdefined\hypertarget\hypertarget{sourcepointer-5}{}\fi
\noindent\textit{Source reference.} Original lemma 0DK0: sources/stacks-chow/tags.txt, source lines 16995--16995. Prerequisite pointer invoked at proof line(s) 3268; it names a standard fact of a sibling chapter that the counted argument uses, and every step of the counted proof built on it is retained verbatim in the delivered body.\par\medskip
\ifdefined\hypertarget\hypertarget{sourcepointer-6}{}\fi
\noindent\textit{Source reference.} Original lemma 07QM: sources/stacks-chow/tags.txt, source lines 9837--9837. Prerequisite pointer invoked at proof line(s) 3274; it names a standard fact of a sibling chapter that the counted argument uses, and every step of the counted proof built on it is retained verbatim in the delivered body.\par\medskip
\ifdefined\hypertarget\hypertarget{sourcepointer-7}{}\fi
\noindent\textit{Source reference.} Original lemma 0DJF: sources/stacks-chow/tags.txt, source lines 16975--16975. Prerequisite pointer invoked at proof line(s) 3278; it names a standard fact of a sibling chapter that the counted argument uses, and every step of the counted proof built on it is retained verbatim in the delivered body.\par\medskip
\ifdefined\hypertarget\hypertarget{sourcepointer-8}{}\fi
\noindent\textit{Source reference.} Original lemma 008H: sources/stacks-chow/tags.txt, source lines 317--317. Prerequisite pointer invoked at proof line(s) 3247; it names a standard fact of a sibling chapter that the counted argument uses, and every step of the counted proof built on it is retained verbatim in the delivered body.\par\medskip
\ifdefined\hypertarget\hypertarget{sourcepointer-9}{}\fi
\noindent\textit{Source reference.} Original lemma 08H5: sources/stacks-chow/tags.txt, source lines 10768--10768. Prerequisite pointer invoked at proof line(s) 3242; it names a standard fact of a sibling chapter that the counted argument uses, and every step of the counted proof built on it is retained verbatim in the delivered body.\par\medskip
\begin{lemma}
\label{lemma-tor-dimension-rel}
Let $f : X \to Y$ be a morphism of schemes. Let $E$ be an object
of $D(\mathcal{O}_X)$. Then
\begin{enumerate}
\item $E$ as an object of $D(f^{-1}\mathcal{O}_Y)$ has tor amplitude in
$[a, b]$ if and only if $\epsilon^*E$ has tor amplitude in $[a, b]$
as an object of $D(f_{small}^{-1}\mathcal{O}_{Y_\etale})$.
\item $E$ locally has finite tor dimension as an object of
$D(f^{-1}\mathcal{O}_Y)$ if and only if $\epsilon^*E$
locally has finite tor dimension as an object of
$D(f_{small}^{-1}\mathcal{O}_{Y_\etale})$.
\end{enumerate}
Here $\epsilon$ is as in (\ref{equation-epsilon}).
\end{lemma}

\begin{proof}
The easy direction in (1) follows from Cohomology on Sites, Lemma
\ref{sites-cohomology-lemma-tor-amplitude-pullback}.
Let $x \in X$ be a point and let $\overline{x}$ be a geometric
point lying over $x$. Let $y = f(x)$ and denote $\overline{y}$
the geometric point of $Y$ coming from $\overline{x}$.
Then $(f^{-1}\mathcal{O}_Y)_x = \mathcal{O}_{Y, y}$
(Sheaves, Lemma \ref{sheaves-lemma-stalk-pullback})
and
$$
(f_{small}^{-1}\mathcal{O}_{Y_\etale})_{\overline{x}} =
\mathcal{O}_{Y_\etale, \overline{y}} =
\mathcal{O}_{Y, y}^{sh}
$$
is the strict henselization
(by \'Etale Cohomology, Lemmas \ref{etale-cohomology-lemma-stalk-pullback}
and \ref{etale-cohomology-lemma-describe-etale-local-ring}).
Since the stalk of $\mathcal{O}_{X_\etale}$ at $X$ is
$\mathcal{O}_{X, x}^{sh}$ we obtain
$$
(\epsilon^*E)_{\overline{x}} =
E_x \otimes_{\mathcal{O}_{X, x}}^\mathbf{L} \mathcal{O}_{X, x}^{sh}
$$
by transitivity of pullbacks. If $\epsilon^*E$ has tor amplitude
in $[a, b]$ as a complex of $f_{small}^{-1}\mathcal{O}_{Y_\etale}$-modules,
then $(\epsilon^*E)_{\overline{x}}$ has tor amplitude in $[a, b]$
as a complex of $\mathcal{O}_{Y, y}^{sh}$-modules
(because taking stalks is a pullback and lemma cited above). By
More on Flatness, Lemma \ref{flat-lemma-tor-amplitude-up-down-henselization}
we find the tor amplitude of
$(\epsilon^*E)_{\overline{x}}$
as a complex of $\mathcal{O}_{Y, y}$-modules is in $[a, b]$.
Since $\mathcal{O}_{X, x} \to \mathcal{O}_{X, x}^{sh}$ is faithfully
flat (More on Algebra, Lemma
\ref{more-algebra-lemma-dumb-properties-henselization}) and since
$(\epsilon^*E)_{\overline{x}} =
E_x \otimes_{\mathcal{O}_{X, x}}^\mathbf{L} \mathcal{O}_{X, x}^{sh}$
we may apply
More on Algebra, Lemma \ref{more-algebra-lemma-no-change-tor-amplitude}
to conclude the tor amplitude of $E_x$ as a complex of
$\mathcal{O}_{Y, y}$-modules is in $[a, b]$.
By Cohomology, Lemma \ref{cohomology-lemma-tor-amplitude-stalk}
we conclude that $E$ as an object of $D(f^{-1}\mathcal{O}_Y)$
has tor amplitude in $[a, b]$. This gives the reverse implication in (1).
Part (2) follows formally from (1).
\end{proof}

\end{document}
